\section{Multiplicity at a logarithmic jet}\label{sec:zero}

We work over $k=\mathbb F_p$. Let $\Sigma\subset k$ be a nonempty set of
distinct finite poles, $b\in k$ a centre, and $(A_a)_{a\in \Sigma}$ constant
strictly upper triangular $N\times N$ matrices in one ordered basis.
Put $T=\prod_{a\in \Sigma\cup\{b\}}(X-a)$ and $d=\deg T-1$. Suppose the
common left kernel of the $A_a$ is generated by a coordinate row $e$.
The polynomial column $w$ in Theorem~\ref{thm:zero} has $ew=1$ and satisfies
\begin{equation}\label{eq:log-jet}
 T w'\equiv\sum_{a\in \Sigma}\frac{T}{X-a}A_aw\pmod{(X-b)^p}.
\end{equation}
At an ordinary centre $b\notin \Sigma$ this jet exists for every initial
column $w_0$ with $ew_0=1$. Indeed, expand $\sum_{a\in\Sigma} A_a/(X-a)$ at $b$
and solve for the coefficients of $w$ through degree $p-1$. Computing
coefficient $j$ requires division by $j$ for $1\le j<p$. The initial
condition and $eA_a=0$ give $ew=1$ for the truncated polynomial.
Multiplication by $T$ raises the order of the remaining differential
defect from at least $p-1$ to at least $p$. At a singular centre
we use \eqref{eq:log-jet} as a hypothesis.

\begin{proof}[Proof of Theorem~\ref{thm:zero}]
On polynomial rows define
\[
 D h=T h'+\sum_{a\in \Sigma}\frac{T}{X-a}h A_a .
\]
If $P$ is a nonzero row of degree at most $A$, then $D^jP$ has entries
of degree at most $A+jd$ by induction. Equation~\eqref{eq:log-jet}
gives $(Dh)w\equiv T(hw)'\pmod{(X-b)^p}$. Since $X-b$ divides $T$,
the operator $T\,d/dX$ preserves $(X-b)^m$ for every $m\ge0$.
Consequently $Pw\equiv0\pmod{(X-b)^m}$ implies
$(D^jP)w\equiv0\pmod{(X-b)^m}$ for all $j$ whenever $m\le p$.

Over $K=k(X)$ take the first $r$ successive independent rows
$P,DP,\ldots,D^{r-1}P$, stopping when the next row is dependent.
Such an $r$ exists with $1\le r\le N$. Write $M$ for their matrix.
Its row span $V$ is stable under $D$, since $D(fh)=Tf'h+fDh$,
and hence also under $\nabla=T^{-1}D$.
Let $I$ be the ordered pivot columns of $M$ and let $B=M_I$.
The rows of $B^{-1}M$ form the reduced echelon basis of $V$.
Its last row $h$ starts with a $1$ in the last pivot column.
In $\nabla h=h'+\sum_{a\in\Sigma} hA_a/(X-a)$, all coordinates through this
column vanish by strict upper triangularity. Every nonzero element
of $V$ starts in a pivot column, so $\nabla h=0$.

All $r$-column minors of $M$ have degree at most
$H_r=rA+d\binom r2$: each determinant uses one entry from each
successive derivative row. By Cramer's rule each coordinate of $h$
is a ratio of such a minor (possibly zero) to $q=\det B\ne0$.
In particular both degrees are bounded by $H_r\le H=NA+d\binom N2<p$.

We now show in increasing order of $j$ that $h_j\in k$.
Horizontality gives
\[
 h_j'=-\sum_{a\in \Sigma}\frac{1}{X-a}\sum_{i<j}h_i(A_a)_{ij}.
\]
By induction all numerators on the right are in $k$.
The formal Laurent derivative at $a$ has zero coefficient at
$(X-a)^{-1}$, since only the derivative of the constant term could
contribute. Taking residues at the distinct $a$ makes every numerator
zero. Hence $h_j'=0$. To see why this implies $h_j\in k$ here, write a
nonzero coordinate in lowest terms as $u/v$. The identity $u'v=uv'$
gives $v\mid v'$, whence $v'=0$ by degree and then $u'=0$.
Both polynomials have degree less than $p$, so coefficient comparison
makes both constant. The zero coordinate is already constant.
This proves the induction, including the empty sum for the first
coordinate. The same residue equations give $hA_a=0$ for every $a$.
The common-kernel hypothesis and the pivot entry $1$ give $h=e$.

The identity $q B^{-1}=\operatorname{adj}(B)$ expresses $qe$ as a
polynomial linear combination of the rows of $M$. If $Pw$ vanished
to order at least $H_r+1$, the propagation of vanishing would apply
at this precision because $H_r+1\le p$. It would give
$(X-b)^{H_r+1}\mid qew=q$, contradicting $0\ne q$ and $\deg q\le H_r$.
Thus $Pw$ has a nonzero coefficient of order at most $H_r\le H$.
\end{proof}

\subsection{Classical finite polylogarithms}
For $1\le i\le n$ put $\ell_i(X)=\sum_{m=1}^{p-1}X^m/m^i$.
The following form of Theorem~\ref{thm:zero} will be used in
Section~\ref{sec:discs}.

\begin{corollary}\label{cor:box}
Let $n\ge1$, $t,A\ge0$, $N=\binom{n+t}{n}$ and
$H=NA+\binom N2$. If $t<p$ and $H<p$, every nonzero
$f\in k[X,y_1,\ldots,y_n]$ of $X$-degree at most $A$ and total
$y$-degree at most $t$ has
\[
 \operatorname{ord}_0 f(X,\ell_1(X),\ldots,\ell_n(X))\le H.
\]
In particular substitution modulo $X^p$ is injective on this degree box.
\end{corollary}

\begin{proof}
Use the $N$ monomials of total degree at most $t$ as coordinates,
ordered by decreasing weight for $\mathrm{wt}(y_i)=i$.
The two derivations
\[
 E=\partial_{y_1},\qquad L=\sum_{i=2}^n y_{i-1}\partial_{y_i}
\]
act on this space by constant strictly upper triangular row matrices:
each lowers weight by one. If $Ef=Lf=0$ on this monomial space, the
identities $[\partial_{y_j},L]=\partial_{y_{j+1}}$ for $1\le j<n$
successively give $\partial_{y_j}f=0$ for every $j$. Each exponent
is below $p$, so $f$ is constant. The common left kernel is precisely
the coordinate line of the monomial $1$.

Take $w$ to be the column obtained by substituting $y_i=\ell_i$ in the
monomials. Differentiating the sums gives
$X(1-X)\ell_1'=X-X^p$ and $X\ell_i'=\ell_{i-1}$ for $i\ge2$.
The product rule now gives \eqref{eq:log-jet} with $b=0$ and residues
$L$ at $0$ and $-E$ at $1$: the connection on monomials is
$L/X+E/(1-X)$. Here $d=1$. Apply the theorem to the row of
coefficients of $f$ to obtain the claimed inequality.
\end{proof}

\subsection{Words and strict multiple polylogarithms}
Take a finite set of distinct words on an alphabet $\Sigma\subset k$,
containing the empty word and closed under deletion of the first letter.
Order the words by decreasing length. The residue for $a$ sends the row
of $av$ to the row of $v$ and kills rows with any other first letter,
including the empty row. These matrices are strictly upper triangular.
In the image under the $a$-matrix, the $v$-coordinate recovers the
original $av$-coefficient whenever $av$ occurs. Thus the common left
kernel is the empty-word line. Multiplying individual residue matrices
by nonzero constants preserves this calculation. Theorem~\ref{thm:zero}
applies to polynomial-coefficient linear forms on these word jets.

For example, with $s_j\ge1$, consider
\[
 F_{s_1,\ldots,s_h}(X)=
 \sum_{p>n_1>\cdots>n_h\ge1}
 \frac{X^{n_1}}{n_1^{s_1}\cdots n_h^{s_h}},\qquad F_{\varnothing}=1 .
\]
Encode this index by $0^{s_1-1}1\cdots0^{s_h-1}1$ and take all suffixes
of the words in question. If $s_1>1$, then
$XF'_{s_1,\mathbf t}=F_{s_1-1,\mathbf t}$. For $s_1=1$ summing the
geometric progression $\sum_{n=m+1}^{p-1}X^n$ above the largest tail
index $m$ (take $m=0$ for the empty tail) gives
\[
 X(1-X)F'_{1,\mathbf t}=XF_{\mathbf t}(X)-X^p F_{\mathbf t}(1).
\]
Modulo $X^p$ the residues are therefore the deletion operators for $0$
and $1$, with a minus sign at $1$. The estimate holds at $b=0$ for
linear forms in distinct words. For a polynomial expression a separate
check of its expansion is needed. For instance $F_1^2-2F_{1,1}$ vanishes
modulo $X^p$ for odd $p$: at each $1\le n<p$ its two terms agree by
$1/(i(n-i))=(1/i+1/(n-i))/n$ summed over $1\le i<n$.

\subsection{Quadratic growth with two fixed poles}
For every $n\ge2$ we construct a system of rank $R$ with
$n+1\le R\le 2n$ and a constant row whose order at an ordinary centre
is at least $n^2-1$. Work first over $\mathbb Q[[X]]$ and define
\[
 J_u f=\int_0^X\frac{f(t)}{1-t}\,dt,\qquad
 J_v f=\int_0^X\frac{f(t)}{1+t}\,dt .
\]
The first $n^2-1$ Taylor coefficients (including the constant) of
\[
 G=\sum_{0\le i,j<n} b_{ij}J_u^iJ_vJ_u^j1
\]
give a rational linear map from dimension $n^2$ to dimension $n^2-1$.
Choose a nonzero kernel element, clear its denominators, and let $k_0$
be the last index of a nonzero row. Introduce the ordered coordinates
$U_0,\ldots,U_{k_0},V_0,\ldots,V_{n-1}$ by the recursions
\[
 V_{n-1}=1,\qquad V_j=J_uV_{j+1}\ (j<n-1),\qquad
 U_i=J_uU_{i+1}+J_v\sum_{j=0}^{n-1}b_{ij}V_{n-1-j},\qquad U_{k_0+1}=0,
\]
where $U_{k_0+1}$ is a boundary value. Expanding recursively gives
$U_0=G$. Differentiation gives a strictly upper triangular system with
poles at $1,-1$ and rank $R=k_0+1+n$. Up to sign the residue at $1$
shifts each of the two chains to the next coordinate. Its left kernel
has basis the terminal rows $e_{U_{k_0}},e_{V_{n-1}}$. The residue at
$-1$ kills $e_{V_{n-1}}$ and takes $e_{U_{k_0}}$ to the nonzero
reversed last array row in the $V$-coordinates. The common kernel is
therefore exactly $k e_{V_{n-1}}$ after reduction at every prime for
which the last array row remains nonzero.

At the ordinary centre $0$ the Taylor coefficients through degree $p-1$
are $p$-integral: the kernels $1/(1\pm X)$ have integral coefficients
and each integration divides the coefficient of degree $j-1$ by $j$.
Reduce this finite jet for an odd prime $p>\max\{n^2,2\binom R2\}$ at
which the last array row survives. Its $U_0$-coordinate has vanishing
order at least $n^2-1$ by construction. Theorem~\ref{thm:zero}, with
$d=2$ and $A=0$, shows that this coordinate is nonzero modulo $X^p$.
As $n$ varies, $R$ is unbounded and $n^2-1\ge R^2/4-1$. Consequently
even two fixed poles and a constant input row admit quadratic growth
in rank. Any improvement to linear growth for the monomial systems of
Corollary~\ref{cor:box} would require a property specific to that family.
